% Generated draft: requires mathematical and source review.
\documentclass[11pt]{article}
\usepackage[margin=25mm]{geometry}
\usepackage{fontspec}
\setmainfont{Latin Modern Roman}
\usepackage{amsmath,amssymb,mathtools,unicode-math}
\setmathfont{Latin Modern Math}
\usepackage{enumitem,needspace,xurl,hyperref}
\setlength{\parindent}{0pt}
\setlength{\parskip}{5pt}
\setlength{\emergencystretch}{3em}
\newcommand{\N}{\mathbb{N}}
\newcommand{\Z}{\mathbb{Z}}
\newcommand{\Q}{\mathbb{Q}}
\newcommand{\R}{\mathbb{R}}
\newcommand{\C}{\mathbb{C}}
\newcommand{\F}{\mathbb{F}}
\newcommand{\A}{\mathbb{A}}
\newcommand{\PP}{\mathbb{P}}
\newcommand{\E}{\mathbb{E}}
\newcommand{\eps}{\varepsilon}
\newcommand{\dd}{\,\mathrm{d}}
\begin{document}
\section*{UnsolvedMath ID: 6000015}
\textbf{Source problem code:} catalogue entry
\subsection{Definitions and mathematical statement}
A Hessian manifold consists of a smooth manifold \(M\), a flat torsion-free affine connection \(D\), and a Riemannian metric \(g\) locally expressible in \(D\)-affine coordinates as \(g_{ij}=\partial_i\partial_j\phi\) for a strictly convex potential. Completeness is taken here as geodesic completeness of \(g\), pending confirmation of the source convention; it is not silently identified with affine completeness of \(D\).

In affine coordinates \((x^i)\), tangent coordinates \((x^i,v^i)\) define complex coordinates \(z^i=x^i+\mathrm i v^i\) on \(TM\). An affine transition \(x'=Ax+b\) induces \(z'=Az+b\), giving the natural integrable complex structure \(J_D\). The first question is
\[
\forall(M,D,g)\text{ complete Hessian},\qquad (TM,J_D)\text{ is Stein}?
\]
Here Stein means, equivalently for the usual finite-dimensional complex-manifold category, admitting a proper holomorphic embedding in some finite-dimensional complex Euclidean space. The choice of flat connection is part of the structure even though suppressed in the source's notation \((M,g)\).

For the specified subquestion, let \(\Omega\subseteq\R^n\) be a nonempty open convex domain containing no complete affine line. Let a discrete affine group \(\Gamma\) preserve \(\Omega\) and act freely and properly discontinuously. Writing \(\gamma(x)=A_\gamma x+b_\gamma\), there is a natural complex identification
\[
T(\Omega/\Gamma)\cong
(\Omega+\mathrm i\R^n)/\Gamma,
\qquad \gamma(z)=A_\gamma z+b_\gamma.
\]
Is this quotient Stein? This is a quotient question, not merely Steinness of the tube before quotienting. The record does not supply a complete invariant Hessian metric for every such quotient or the precise implication intended by in particular.

The supplied research summary reports the Sasaki lift as Kähler and mentions recent lifting, positivity, and splitting work, but says that no published general Steinness theorem was verified. Kählerness alone is not substituted for the requested Stein conclusion; the alleged classical convex-domain subcase is described as plausible rather than verified.
\subsection{Short English statement}
Is the natural complex tangent bundle of every complete Hessian manifold Stein? In particular, is the tube quotient \((\Omega+\mathrm i\R^n)/\Gamma\) Stein under the stated convexity and affine-action assumptions?
\subsection{Sources}
\begin{itemize}
\item Local UnsolvedMath registry, permanent ID \(6000015\), exact source code \texttt{catalogue entry}, ``Stein Tangent Bundles of Complete Hessian Manifolds.'' Proposer Satoru Shimizu, 1998.
\item Furuhata, Matsuzoe and Urakawa, \emph{Open Problems in Affine Differential Geometry and Related Topics} (1998), item 5(a), recorded scanned PDF pages 1--2; \url{https://www.jstage.jst.go.jp/article/iis/4/2/4_2_125/_pdf/-char/en}. Recorded access 2026-07-29; manual OCR from page images.
\item Supplied research summary mentions arXiv:2509.01176 for Sasaki-lift applications involving positive flat/holomorphic line bundles and fibration/splitting results for compact Hessian manifolds. No title, authors, or theorem implying the requested Stein conclusion are supplied.
\end{itemize}
\end{document}
